
\chapter{Fundamentação Teórica}
\label{chap:fundamentacao}

% Apresente os conceitos que fundamentam o desenvolvimento
% do DoseCerta. Explique como eles se relacionam ao problema
% abordado e às escolhas técnicas do projeto.

\section{Polifarmácia e organização do uso de medicamentos}
\label{sec:polifarmacia}

% Defina polifarmácia com base em fontes acadêmicas.
% Aborde os desafios relacionados à organização de múltiplos
% medicamentos, doses e horários.

\section{Tecnologias de apoio à administração de medicamentos}
\label{sec:tecnologias_medicamentos}

% Apresente soluções tecnológicas utilizadas para apoiar
% a organização e o acompanhamento de medicamentos.
% Diferencie lembretes digitais de dispositivos dispensadores.

\section{Sistemas embarcados}
\label{sec:sistemas_embarcados}

% Explique o conceito de sistemas embarcados e sua aplicação
% em dispositivos dedicados.
% Relacione os conceitos à plataforma utilizada no projeto.

\subsection{Plataforma de processamento}
\label{subsec:plataforma_processamento}

% Apresente a plataforma computacional escolhida,
% suas características e os motivos da escolha.

\section{Sensores e sistemas de acionamento}
\label{sec:sensores_acionamento}

% Apresente os princípios dos sensores e dos mecanismos
% utilizados no dispositivo.
% Explique os conceitos relevantes para o projeto mecânico
% e eletrônico.

\section{Visão computacional}
\label{sec:visao_computacional}

% Introduza os conceitos de visão computacional pertinentes
% ao projeto, incluindo aquisição e processamento de imagens,
% caso sejam utilizados.

\section{Aplicações web e comunicação entre sistemas}
\label{sec:aplicacoes_web}

% Apresente os conceitos necessários para compreender
% o backend, a interface do usuário e a comunicação
% entre os componentes do sistema.

\subsection{Interfaces e experiência do usuário}
\label{subsec:interface_usuario}

% Aborde princípios de interface e usabilidade
% relevantes para o público-alvo.

\subsection{Serviços de backend e APIs}
\label{subsec:backend_apis}

% Explique os conceitos de backend, APIs e integração
% entre os componentes da aplicação.

\section{Inteligência artificial e agentes conversacionais}
\label{sec:inteligencia_artificial}

% Apresente os conceitos de inteligência artificial
% utilizados no projeto.

\subsection{Modelos de linguagem}
\label{subsec:modelos_linguagem}

% Explique os conceitos relevantes dos modelos de linguagem
% e suas possibilidades de aplicação.

\subsection{Agentes de inteligência artificial}
\label{subsec:agentes_ia}

% Descreva o conceito de agente e sua função no sistema,
% delimitando suas capacidades e restrições.

\section{Considerações sobre segurança e confiabilidade}
\label{sec:seguranca_confiabilidade}

% Apresente conceitos relacionados à segurança do dispositivo,
% à confiabilidade do registro de informações e aos limites
% de atuação de uma solução de apoio à organização
% de medicamentos.
